% TEMPLATE-MILC-v3.tex - plantilla maestra UNICA de las guias MILC v3.
%
% La usan las cuatro familias de guias, cada una con su driver:
%   · Grados 6-11          -> scripts/build-guias-g11.py
%   · Bebras               -> scripts/build-guias-bebras.py
%   · Semillero            -> scripts/build-guias-semillero.py
%   · Territorio interior  -> scripts/build-guias-territorio-interior.py
%
% Antes habia tres copias casi identicas (~400 lineas iguales, 6 distintas):
% cada mejora habia que aplicarla tres veces, y en la practica no se hacia
% (el motor de recursos vivia solo en dos de ellas). Lo que varia entre
% familias esta parametrizado en 6 placeholders que cada driver llena
% (se nombran aqui SIN su sintaxis real: el builder reemplaza por texto plano,
% tambien dentro de los comentarios, y un valor multilinea rompe el preambulo):
%
%   HEADER_IZQ        encabezado izquierdo de pagina
%   HEADER_DER        encabezado derecho de pagina
%   PORTADA_ETIQUETA  rotulo sobre el numero grande (GRADO/MOMENTO/MODULO)
%   PORTADA_SERIE     linea de serie de la portada
%   PORTADA_ID        identificador de la guia en la portada
%   PORTADA_CIRCULO   el circulito de grado (°), o vacio si no aplica
%   PDF_TITULO        titulo en texto plano (sin \textbf ni comillas TeX) para
%                     los metadatos del PDF (/Title); lo produce lib_xelatex.texto_pdf
%
% Composicion (auditoria 2026-09-03): las bandas de estacion piden espacio
% (\Needspace*) y prohiben el salto justo despues (\nopagebreak), asi no quedan
% huerfanas al pie; las cajas largas (softbox, drawbox, infoband, puentebox) son
% `breakable`, asi una caja de media pagina ya no arrastra todo a la siguiente
% dejando la anterior al 40 %. Todo texto sobre mostaza o verde va en milcNegro
% (blanco sobre #E5B400 da 1,93:1 y sobre #58A923 2,95:1; ninguno pasa WCAG AA).
% El resto de claves las documenta content/guias/_SCHEMA.md. El script valida
% que no queden placeholders sin reemplazar antes de escribir el archivo final.

% Etiquetado PDF (latex-lab): /Lang, estructura y texto alternativo de las
% figuras, para lectores de pantalla. Debe ir ANTES de \documentclass.
\DocumentMetadata{lang=es-CO,tagging=on}
\documentclass[11pt,letterpaper]{article}
% headheight=24pt: la cabecera derecha lleva dos lineas (identidad de la guia
% y estacion actual). headsep se acorta en la misma medida para que el bloque
% de texto quede exactamente donde estaba con headheight=12pt.
\usepackage[margin=2.0cm,headheight=24pt,headsep=13pt]{geometry}
\usepackage[table]{xcolor}
\usepackage{fontspec}
\usepackage[spanish]{babel}
\usepackage{tikz}
% Librerías TikZ para los diagramas de las guías (bloque `recursos:`):
%  · babel  → babel en español vuelve activos caracteres como «>», y TikZ los usa
%             en sus opciones (>=stealth, ->). Sin esta librería, cualquier
%             diagrama con flechas revienta con «\language@active@arg> has an
%             extra }». Debe ir ANTES de que se dibuje nada.
%  · positioning        → sintaxis «below left=2pt and 3pt».
%  · decorations.pathreplacing → llaves ({brace}) para señalar rangos.
\usetikzlibrary{babel,positioning,decorations.pathreplacing}
\usepackage{graphicx}
% Circuitos: el trabajo de electrónica se hace hoy con micro:bit y sus sensores
% integrados (MakeCode, sin cableado), así que no se necesitan esquemáticos.
% Si algún día entran montajes con protoboard, basta descomentar esta línea:
% los diagramas .tex/.tikz declarados en `recursos:` ya se incrustan solos.
% \usepackage{circuitikz}
\usepackage[most]{tcolorbox}
\usepackage{tabularx}
\usepackage{array}
\usepackage{fancyhdr}
\usepackage{titlesec}
\usepackage[document]{ragged2e}  % bandera a la derecha en todo el cuerpo
\usepackage{enumitem}
\usepackage{microtype}
\usepackage{etoolbox}
\usepackage{needspace}
% hyperref al final del preambulo: metadatos del PDF (titulo, autor, idioma,
% familia), marcadores por estacion (\pdfbookmark) y enlaces sin recuadro.
\usepackage[hidelinks,
  pdftitle={Modelo de negocio — el lienzo de tu proyecto},
  pdfauthor={PhD. Álvaro Cárdenas Orozco},
  pdfsubject={Tecnología e Informática · Grado 11 · Guía 23 · MILC v3},
  pdflang={es-CO},
  bookmarksopen=true,bookmarksnumbered=false]{hyperref}

% Helvetica Neue no trae la flecha U+2192, asi que cada «→» escrito literalmente
% en el YAML se compilaba a NADA, en silencio (462 flechas en 61 guias: rutas del
% tipo «paleta → tipografia → lockup» salian rotas). Mapeamos el caracter a su
% version matematica, que si tiene glifo. Asi el autor puede seguir escribiendo
% «→» comodamente en el YAML.
\usepackage{newunicodechar}
\newunicodechar{→}{\ensuremath{\rightarrow}}
% Mismo problema con los simbolos de marca y vinetas: el «✓» de «una palomita ✓»
% salia como cuadrito vacio en el PDF de todas las guias que lo usaban.
\usepackage{pifont}
\newunicodechar{✓}{\ding{51}}
\newunicodechar{✗}{\ding{55}}
\newunicodechar{✔}{\ding{52}}
\newunicodechar{•}{\textbullet}
\newunicodechar{≥}{\ensuremath{\geq}}
\newunicodechar{≤}{\ensuremath{\leq}}

\setmainfont{Helvetica Neue}
\setsansfont{Helvetica Neue}
\setlength{\parindent}{0pt}
\setlength{\parskip}{6pt}
% Sin particion de palabras: en 6.º-7.º una palabra cortada por guion al final
% de linea («Comuni-cacion») frena la lectura mucho mas que una linea corta.
\hyphenpenalty=10000
\exhyphenpenalty=10000
\pretolerance=2000
\tolerance=3000
\setlength{\emergencystretch}{3em}
\linespread{1.10}
\renewcommand{\arraystretch}{1.25}
\setlist{leftmargin=5mm,itemsep=1mm,topsep=1mm}
\newcolumntype{Y}{>{\RaggedRight\arraybackslash}X}

\definecolor{milcMagenta}{HTML}{E6007E}
\definecolor{milcVino}{HTML}{5A0038}
\definecolor{milcTurquesa}{HTML}{0093A5}
\definecolor{milcVerde}{HTML}{58A923}
\definecolor{milcMostaza}{HTML}{E5B400}
\definecolor{milcOcre}{HTML}{B86600}
\definecolor{milcNegro}{HTML}{241816}
\definecolor{milcGris}{HTML}{F7F7F4}
\definecolor{milcLinea}{HTML}{E8E2DD}
\definecolor{milcGradePrimary}{HTML}{0066FF}
\definecolor{milcGradeDark}{HTML}{003D99}
\definecolor{milcGradeSoft}{HTML}{D6E8FF}
\definecolor{milcGradeAccent}{HTML}{CFE7FF}
\definecolor{milcGradeText}{HTML}{FFFFFF}
\color{milcNegro}

\pagestyle{fancy}
\fancyhf{}
% Cabecera de dos lineas: identidad de la guia arriba y, a la derecha y en
% gris, la estacion en curso (\markright desde \milcestacion). Antes de la
% primera estacion la segunda linea va vacia; el \strut mantiene la altura
% para que la linea de arriba no baile entre paginas.
\lhead{\small Tecnología e Informática · Grado 11\\[-1pt]{\footnotesize\strut}}
\rhead{\small Guía 23 · MILC v3\\[-1pt]{\footnotesize\strut\color{milcNegro!65}\rightmark}}
\cfoot{\small\thepage}
\renewcommand{\headrulewidth}{0pt}

% Sobre mostaza (#E5B400) y verde (#58A923) el texto va en milcNegro; sobre el
% resto de la paleta, en blanco. Se usa dentro de fonttitle/fontupper, donde un
% \color posterior manda sobre coltitle.
\newcommand{\milctintasobre}[1]{%
  \ifboolexpr{test{\ifstrequal{#1}{milcMostaza}} or test{\ifstrequal{#1}{milcVerde}}}%
    {\color{milcNegro}}{\color{white}}}

\newtcolorbox{titlebox}[1]{
  enhanced,colback=#1,colframe=#1,arc=7mm,boxrule=0pt,
  left=7mm,right=7mm,top=3mm,bottom=3mm,
  before skip=8pt,after skip=8pt,
  fontupper=\sffamily\bfseries\Large\color{white}
}

\newtcolorbox{softbox}[3]{
  enhanced,breakable,pad at break=3mm,
  colback=#2,colframe=#1,borderline west={4pt}{0pt}{#1},
  arc=2mm,boxrule=.4pt,left=4mm,right=4mm,top=3mm,bottom=3mm,
  before skip=6pt,after skip=8pt,title={#3},coltitle=white,
  colbacktitle=#1,fonttitle=\sffamily\bfseries\milctintasobre{#1},fontupper=\RaggedRight,
  attach boxed title to top left={xshift=3mm,yshift=-2mm},
  boxed title style={arc=2mm, boxrule=0pt}
}

\newtcolorbox{drawbox}[2]{
  enhanced,breakable,pad at break=4mm,
  colback=white,colframe=#1,arc=4mm,boxrule=1pt,
  left=5mm,right=5mm,top=4mm,bottom=4mm,before skip=8pt,after skip=8pt,
  title={#2},coltitle=white,colbacktitle=#1,fonttitle=\sffamily\bfseries\milctintasobre{#1},
  fontupper=\RaggedRight,
  attach boxed title to top left={xshift=4mm,yshift=-2mm},
  boxed title style={arc=3mm, boxrule=0pt}
}

\newtcolorbox{infoband}[1]{
  enhanced,breakable,pad at break=2mm,colback=#1!8,colframe=#1!50,arc=2mm,boxrule=.5pt,
  left=4mm,right=4mm,top=2mm,bottom=2mm,before skip=4pt,after skip=4pt,
  fontupper=\small\RaggedRight
}

% Puente narrativo entre secciones (contrato editorial Conéctate).
% Da continuidad: "Acabas de X. Ahora vas a Y porque Z."
% Visualmente: banda izquierda en color del grado, texto en itálica pequeña.
\newtcolorbox{puentebox}{
  enhanced,breakable,pad at break=2mm,colback=milcGradeSoft,colframe=milcGradeDark,
  borderline west={3pt}{0pt}{milcGradeDark},
  arc=0mm,boxrule=0pt,left=5mm,right=4mm,top=2.5mm,bottom=2.5mm,
  before skip=4pt,after skip=8pt,
  fontupper=\itshape\small\color{milcNegro}\RaggedRight
}
% La flecha va en modo matematico a proposito: Helvetica Neue NO tiene el glifo
% U+2192, asi que \textrightarrow se compilaba a NADA («Missing character: There
% is no → in font Helvetica Neue») y el puente salia sin flecha. En modo
% matematico la flecha viene de la fuente matematica y si se dibuja.
\newcommand{\puente}[1]{%
  \begin{puentebox}%
    \textcolor{milcGradeDark}{$\rightarrow$}\hspace{2mm}#1%
  \end{puentebox}%
}

\newcommand{\linead}[1]{\noindent\color{milcLinea}\rule{#1}{0.5pt}\color{milcNegro}}
\newcommand{\respuesta}[1]{\par\vspace{2mm}\foreach \n in {1,...,#1}{\noindent\color{milcLinea}\rule{\linewidth}{0.5pt}\par\vspace{5mm}}\color{milcNegro}}
\newcommand{\checkbox}{\textcolor{milcGradeDark}{\fbox{\phantom{x}}}\hspace{5pt}}

% ─── Listas aireadas para las guías socioemocionales ────────────────────
% Componer las verificaciones como «\checkbox texto\\» deja el texto que salta
% de línea debajo del cuadrito y apelmaza el bloque. Estos entornos alinean la
% continuación y airean los ítems. Los usa Territorio Interior; cualquier
% familia puede hacerlo.
\newlist{checklista}{itemize}{1}
\setlist[checklista]{
  label=\textcolor{milcGradeDark}{\fbox{\phantom{x}}},
  leftmargin=9mm, labelsep=3.2mm, itemsep=2.4mm,
  topsep=2.5mm, parsep=0pt, partopsep=0pt, before=\RaggedRight
}
\newlist{pasoslista}{enumerate}{1}
\setlist[pasoslista]{
  label=\textcolor{milcGradeDark}{\sffamily\bfseries\arabic*.},
  leftmargin=8mm, labelsep=3mm, itemsep=2.8mm,
  topsep=1mm, parsep=0pt, partopsep=0pt, before=\RaggedRight
}

\newcommand{\phasecard}[4]{
\begin{tcolorbox}[enhanced,colback=#3,colframe=#2,arc=3mm,boxrule=.5pt,height=3.05cm,valign=center,halign=center,left=1.5mm,right=1.5mm,top=2mm,bottom=2mm]
{\sffamily\bfseries\fontsize{22}{24}\selectfont\textcolor{#2}{#1}}\par
{\sffamily\bfseries\small #4}
\end{tcolorbox}
}

% ── Piezas del rediseno 2026 (legibilidad 6.º-11.º) ───────────────────
% Banda de estacion: reemplaza el titlebox macizo. Titulo a la izquierda,
% dato practico (tiempo/modalidad) a la derecha, en una sola linea.
% Nunca huerfana: pide 9 lineas libres (o salta de pagina rellenando con
% \vfill) y prohibe el salto justo despues de la banda. Nueve y no seis:
% la caja partible que sigue necesita ~80pt (titulo adosado, relleno y dos
% lineas) para arrancar en la misma pagina; con menos, tcolorbox la manda
% entera a la siguiente y la banda se queda sola. El penalty 10000 va
% en `after`, ANTES del espacio posterior: un `after skip` (aunque sea 0pt)
% es un \vskip tras la caja, y ese glue es un punto de corte legal que un
% \nopagebreak escrito despues de \end{tcolorbox} ya no cierra. Cada banda
% deja marcador PDF y actualiza la cabecera derecha.
\newcounter{milcestacion}
\newcommand{\milcestacion}[2]{%
\Needspace*{9\baselineskip}%
\stepcounter{milcestacion}%
\pdfbookmark[1]{#2}{milcestacion\themilcestacion}%
\markright{#2}%
\begin{tcolorbox}[enhanced,colback=#1,colframe=#1,arc=2mm,boxrule=0pt,
  left=5mm,right=5mm,top=2.6mm,bottom=2.6mm,before skip=11pt,
  after={\par\nopagebreak\vskip7pt}]
{\sffamily\bfseries\fontsize{15}{18}\selectfont\milctintasobre{#1}#2}
\end{tcolorbox}}

% Tarjeta de paso numerada: sustituye las tablas «Pilar / Aplicacion», que
% eran formato de consulta docente, no de estudiante. El numero va en un
% circulo sobre el borde para que la secuencia se lea de un vistazo.
\newcommand{\milcpaso}[3]{%
\Needspace{4\baselineskip}%
\begin{tcolorbox}[enhanced,colback=white,colframe=milcLinea,arc=1.5mm,boxrule=.9pt,
  left=14mm,right=4mm,top=3mm,bottom=3mm,before skip=4pt,after skip=4pt,
  fontupper=\RaggedRight,
  overlay={\node[anchor=north west,circle,fill=#2,text=white,inner sep=0pt,
    minimum size=7.4mm,font=\sffamily\bfseries\small]
    at ([xshift=3.6mm,yshift=-3.2mm]frame.north west) {#1};}]
#3
\end{tcolorbox}}

% Casilla marcable de tamano util con lapiz (la anterior era \fbox{\phantom{x}},
% de alto variable segun el texto que la seguia).
\newcommand{\milcmarca}{\framebox[3.8mm]{\rule{0pt}{3.4mm}}}

% Fila marcable de lista suelta (checks de la estacion 1).
\newcommand{\milccheckrow}[1]{%
\par\vspace{1.8mm}\noindent
\makebox[6.4mm][l]{\raisebox{-.55mm}{\textcolor{milcNegro}{\milcmarca}}}%
\parbox[t]{\dimexpr\linewidth-6.4mm\relax}{\RaggedRight #1}\par}

% Celda marcable dentro de la rubrica (tabla de autoevaluacion). Misma
% sangria francesa, pero pensada para vivir dentro de una columna X.
\newcommand{\milcopcion}[1]{%
\makebox[5.2mm][l]{\raisebox{-.55mm}{\textcolor{milcNegro}{\milcmarca}}}%
\parbox[t]{\dimexpr\linewidth-5.2mm\relax}{\RaggedRight #1}}

% ── Recursos visuales de la guía ──────────────────────────────────────
% Declarados en el bloque `recursos:` del YAML y emitidos por el builder.
% \guiaFigura   → imagen rasterizada o PDF (foto, carta celeste, montaje).
% \guiaDiagrama → fuente TikZ/circuitikz (.tex) incrustada como vector:
%                 es la vía para circuitos y esquemáticos editables.
% [#1] = texto alternativo (alt) · #2 = ruta relativa al .tex · #3 = pie
% El alt llega al PDF etiquetado (/Alt) y lo lee el lector de pantalla.
\newcommand{\guiaFigura}[3][]{%
\begin{center}
\includegraphics[width=0.88\linewidth,height=0.38\textheight,keepaspectratio,alt={#1}]{#2}\\[1.5mm]
{\footnotesize\itshape\textcolor{milcNegro!75}{#3}}
\end{center}
}

\newcommand{\guiaDiagrama}[2]{%
\begin{center}
\input{#1}\\[1.5mm]
{\footnotesize\itshape\textcolor{milcNegro!75}{#2}}
\end{center}
}

\begin{document}
\thispagestyle{empty}

% ── Portada ───────────────────────────────────────────────────────────
% Antes: un rectangulo de color a pagina completa. Se veia bien en pantalla,
% pero en la fotocopiadora del colegio se comia el toner y salia gris sucio.
% Ahora la banda ocupa el tercio superior y el resto va en blanco.
\begin{tikzpicture}[remember picture,overlay]
  \path[fill=milcGradePrimary]
    (current page.north west) rectangle ([yshift=-10.6cm]current page.north east);

  \node[anchor=north west,text=milcGradeText] at ([xshift=2.0cm,yshift=-1.55cm]current page.north west)
    {\sffamily\bfseries\fontsize{12}{14}\selectfont GRADO};
  \node[anchor=north west,text=milcGradeText,opacity=.85] at ([xshift=2.0cm,yshift=-2.35cm]current page.north west)
    {\sffamily\bfseries\fontsize{8.5}{10}\selectfont SERIE GUÍAS · TECNOLOGÍA E INFORMÁTICA};
  \node[anchor=north east,text=milcGradeText,opacity=.85,align=right] at ([xshift=-2.0cm,yshift=-1.55cm]current page.north east)
    {\sffamily\bfseries\fontsize{9}{12}\selectfont I.E. Sor María Juliana\\Cartago, Valle del Cauca};

  \node[anchor=west,text=milcGradeText] (gradonum) at ([xshift=1.85cm,yshift=-7.1cm]current page.north west)
    {\sffamily\bfseries\fontsize{112}{112}\selectfont 11};
  % El circulito de grado (°) solo aplica a las guias de grado («11°»); en
  % Bebras y Semillero el numero grande es un momento/modulo, no un grado.
  % Se posiciona relativo al nodo `gradonum`, no en coordenadas absolutas.
  \draw[milcGradeText,line width=1.6pt]
    ([xshift=2.0mm,yshift=-4.5mm]gradonum.north east) circle (.15cm);

  \node[anchor=west,text=milcGradeText,text width=11.4cm,align=left] at ([xshift=1.5cm,yshift=0.5cm]gradonum.east)
    {
     {\sffamily\bfseries\fontsize{9}{11}\selectfont\MakeUppercase{Período 3 · Proyecto final emprendedor}}\\[3mm]
     {\sffamily\bfseries\fontsize{21}{25}\selectfont\setlength{\baselineskip}{27pt}Modelo de negocio ---\\[4pt]el lienzo\\[4pt]de tu proyecto}\\[3.5mm]
     {\sffamily\bfseries\fontsize{15}{18}\selectfont Guía 23-11-TIC}
    };
\end{tikzpicture}

\vspace*{9.4cm}
\vfill

{\sffamily\bfseries\fontsize{8}{10}\selectfont\textcolor{milcGradeDark}{LO QUE TE LLEVAS HOY}}

\begin{tcolorbox}[enhanced,colback=white,colframe=milcNegro,arc=2mm,boxrule=1.6pt,
  left=5mm,right=5mm,top=4mm,bottom=4mm,before skip=3pt,after skip=12pt,
  fontupper=\RaggedRight\fontsize{11}{15.5}\selectfont]
Business Model Canvas completo, en una hoja, con los 9 bloques llenos y construidos a partir de los hallazgos de las 5 entrevistas de la sesión 2 (segmento, propuesta de valor, canales, relación con clientes, fuentes de ingreso, recursos clave, actividades clave, socios clave, estructura de costos) + tabla de 3 supuestos críticos por verificar antes de prototipar
\end{tcolorbox}

\begin{tcolorbox}[enhanced,colback=milcGris,colframe=milcGris,arc=2mm,boxrule=0pt,
  left=5mm,right=5mm,top=3.5mm,bottom=3.5mm,after skip=12pt]
{\sffamily\bfseries\fontsize{8}{10}\selectfont\textcolor{milcNegro!55}{ESTA GUÍA ES DE}}\\[3mm]
\begin{tabularx}{\linewidth}{X p{3.2cm} p{3.2cm}}
\linead{\linewidth} & \linead{3.0cm} & \linead{3.0cm}\\[-1mm]
{\fontsize{8.5}{10}\selectfont\textcolor{milcNegro!55}{Nombre completo}} &
{\fontsize{8.5}{10}\selectfont\textcolor{milcNegro!55}{Grupo}} &
{\fontsize{8.5}{10}\selectfont\textcolor{milcNegro!55}{Fecha}}\\
\end{tabularx}
\end{tcolorbox}

\vfill

\noindent\textcolor{milcLinea}{\rule{\linewidth}{1pt}}\\[2mm]
{\sffamily\bfseries\fontsize{9.5}{12}\selectfont Autor: Dr. Álvaro Cárdenas Orozco}\\[1mm]
{\sffamily\fontsize{8.5}{11}\selectfont\textcolor{milcNegro!55}{Metodología MILC · apertura ancestral · pensamiento computacional · triángulo Dussel–estoicismo–Floridi}}

\newpage

\section*{Apertura: Modelo de negocio --- el lienzo de tu proyecto}

\begin{softbox}{milcGradeDark}{milcGradeSoft}{Saber ancestral}
En la galería de mercado de Cartago y en las plazas del Pacífico, el oficio más antiguo del mundo opera todos los sábados sin que nadie lo llame \emph{``modelo de negocio''}: el \textbf{trueque}. La señora trae racimo de plátano y se va con pescado; el pescador llega con corvinas y vuelve con queso; el quesero baja con cuajada y sube con frutas. Ese intercambio ancestral tiene una arquitectura silenciosa que cualquier emprendedor moderno reconocería: hay un \textbf{segmento} (quién intercambia con quién), una \textbf{propuesta de valor} (qué da cada uno), un \textbf{canal} (la plaza, el día, la hora), una \textbf{relación} (la confianza acumulada de años), una \textbf{moneda implícita} (volumen, calidad, frescura) y unos \textbf{recursos clave} (terreno, semilla, embarcación, conocimiento). El \textbf{Business Model Canvas} moderno no inventa nada nuevo: \textbf{ordena en 9 casillas lo que el trueque hacía intuitivamente}. La diferencia es la escala y la urgencia de hacerlo explícito.
\end{softbox}

\begin{softbox}{milcTurquesa}{milcGris}{Saber contemporáneo}
El \textbf{Business Model Canvas} (Alex Osterwalder, 2010) es la herramienta visual más usada del mundo para diseñar un modelo de negocio en una sola hoja. Sus 9 bloques se organizan en 4 zonas: \textbf{Frente} (lo que el cliente ve y vive): segmento, propuesta de valor, canales, relación. \textbf{Atrás} (lo que la empresa hace para entregar): recursos clave, actividades clave, socios clave. \textbf{Base} (la economía): estructura de costos, fuentes de ingreso. La regla profesional: \textbf{el Canvas no es planeación, es hipótesis}. Cada uno de los 9 bloques es una afirmación que aún debes verificar con el mundo. Los emprendedores novatos llenan el Canvas como si fuera un examen; los profesionales lo llenan sabiendo que \textbf{2 o 3 bloques contienen los supuestos más críticos} y la siguiente sesión los pondrá a prueba. Sin Canvas no hay claridad estratégica; pero el Canvas sin validación es solo un cuadro bonito.
\end{softbox}

\begin{softbox}{milcMostaza}{milcGris}{Pregunta-puente}
\textit{¿Qué arquitectura silenciosa del trueque ancestral aún sostiene los modelos de negocio modernos, y por qué el Canvas se llama \emph{lienzo} (canvas) y no \emph{plan}? ¿Qué pasa cuando un estudiante llena el Canvas sin haber hecho las entrevistas de la sesión 2?}
\end{softbox}

\begin{softbox}{milcVerde}{milcGris}{Saber y saber hacer}
Al terminar podrás: (1) \textbf{analizar} los 9 bloques del Canvas y cómo se relacionan entre sí; (2) \textbf{explicar} la diferencia entre Canvas (hipótesis visual) y plan de negocio tradicional (documento de 50 páginas); (3) \textbf{crear} tu Canvas en una hoja con los 9 bloques llenos a partir de la evidencia de las entrevistas de la sesión 2; (4) \textbf{evaluar} cuáles son los 3 supuestos más críticos de tu modelo (los que, si fallan, todo se cae) para llevarlos a verificación.
\end{softbox}



\small
\begin{tabularx}{\textwidth}{>{\bfseries}p{3.0cm}Y}
\rowcolor{milcGradePrimary!12}Autor & Dr. Álvaro Cárdenas Orozco\\
DBA & Diseño, implementación y sustentación de un proyecto de impacto local (MEN, T\&I)\\
Referentes & Dussel (analéctica · sur global) · Lean Startup · Floridi · Estoicismo\\
Producto final & Business Model Canvas completo, en una hoja, con los 9 bloques llenos y construidos a partir de los hallazgos de las 5 entrevistas de la sesión 2 (segmento, propuesta de valor, canales, relación con clientes, fuentes de ingreso, recursos clave, actividades clave, socios clave, estructura de costos) + tabla de 3 supuestos críticos por verificar antes de prototipar\\
\end{tabularx}
\normalsize

\puente{Antes de empezar, mira el viaje completo. En la sesión 1 elegiste un problema; en la sesión 2 lo validaste con 5 entrevistas. Hoy diseñas la \textbf{arquitectura económica} de tu solución en 9 casillas. Las próximas sesiones tomarán este Canvas como mapa para construir prototipo (S4), marca (S5-S6), finanzas (S7), pitch (S8) y sustentación (S9-S10).}

\milcestacion{milcGradeDark}{Ruta MILC: así vamos a aprender}
\begin{tabularx}{\textwidth}{>{\centering\arraybackslash}X>{\centering\arraybackslash}X>{\centering\arraybackslash}X>{\centering\arraybackslash}X}
\phasecard{1}{milcVerde}{milcGris}{Escucha\\[-1mm]\small Observo, escucho y pregunto.} &
\phasecard{2}{milcTurquesa}{milcGris}{Entender\\[-1mm]\small Sistematizo y ordeno.} &
\phasecard{3}{milcMagenta}{milcGris}{Hacer\\[-1mm]\small Construyo y aplico.} &
\phasecard{4}{milcVino}{milcGris}{Cierre\\[-1mm]\small Evalúo y reflexiono.}
\end{tabularx}


\puente{Antes de teorizar, vas a hacer una \emph{lectura honesta} de tu síntesis de la sesión 2. Vas a extraer 3 cosas: quién es exactamente el cliente, qué dolor preciso le quitas y qué pista de canal apareció en las entrevistas. Esos 3 elementos ya tienen 3 bloques del Canvas pre-llenados con evidencia real.}

\milcestacion{milcVerde}{Estación 1 de 3 · Escucha}

\begin{softbox}{milcVerde}{milcGris}{Escena para escuchar}
\textbf{Actividad 1 · ANALIZA --- Pre-llenado del Canvas desde mis entrevistas} (15 min · individual). Toma la síntesis de la sesión 2 y extrae con marcador 3 elementos para pre-llenar el Canvas: (1) \textbf{Segmento concreto}: ¿qué tienen en común los 5 entrevistados que sufren el problema? (oficio, edad, ubicación, contexto). (2) \textbf{Dolor verificado}: la frase textual que más se repitió en las entrevistas describiendo el dolor. (3) \textbf{Pista de canal}: ¿dónde están estas personas?, ¿cómo se enteran de cosas nuevas? (esto ya apareció en las entrevistas si preguntaste bien). No llenes todavía los 9 bloques: solo estos 3 elementos base.
\end{softbox}

{\sffamily\bfseries\fontsize{8}{10}\selectfont\textcolor{milcNegro!55}{MARCA CUANDO LO HAYAS HECHO}}
\milccheckrow{Definí el segmento con 3-4 rasgos concretos (no \emph{``los jóvenes''}).}
\milccheckrow{Tengo una frase textual del entrevistado que describe el dolor.}
\milccheckrow{Tengo una pista real de canal (no inventada).}

\begin{infoband}{milcVerde}


\textbf{Lo que va al cuaderno (Actividad 1 · ANALIZA).} Título: «Actividad 1 --- Pre-llenado del Canvas». \textbf{Formato:} ficha con 3 campos llenos (segmento, dolor verificado con cita, pista de canal). \textbf{Sabes que terminaste cuando} estos 3 elementos vienen de evidencia, no de tu cabeza.
\end{infoband}

\puente{Ya tienes el material crudo. Ahora vas a entender la \textbf{anatomía del Canvas}: los 9 bloques, su orden de llenado correcto (no se llena de izquierda a derecha) y los 5 errores típicos (Canvas hecho desde la idea de la solución, segmento difuso, propuesta de valor genérica, ingresos sin números, costos olvidados).}

\milcestacion{milcTurquesa}{Estación 2 de 3 · Entender}

\begin{softbox}{milcTurquesa}{milcGris}{Lo que vas a entender}
El Canvas tiene 9 bloques, pero no se llena leyendo de izquierda a derecha. Tiene un \textbf{orden profesional} de llenado y unos errores recurrentes que delatan al novato. Vamos a descomponer la herramienta para que tu lienzo de hoy resista la lectura de un asesor.
\end{softbox}

\milcpaso{1}{milcTurquesa}{Los 9 bloques se distribuyen en 4 zonas: (1) \textbf{Zona del cliente} (frente): \emph{Segmento de cliente} (¿quién?), \emph{Propuesta de valor} (¿qué dolor le quitas o qué ganancia le das?), \emph{Canales} (¿cómo lo alcanzas?), \emph{Relación con clientes} (¿cómo interactúan tú y él?). (2) \textbf{Zona económica} (base): \emph{Fuentes de ingreso} (¿por qué te paga?, ¿cuánto?, ¿cómo?), \emph{Estructura de costos} (¿qué te cuesta entregar?). (3) \textbf{Zona operativa} (atrás): \emph{Recursos clave} (¿qué necesitas tener?), \emph{Actividades clave} (¿qué haces tú?), \emph{Socios clave} (¿quién hace lo que tú no haces?).}
\milcpaso{2}{milcTurquesa}{El \textbf{orden profesional de llenado} es: (1) Segmento → (2) Propuesta de valor → (3) Canales → (4) Relación → (5) Fuentes de ingreso → (6) Recursos clave → (7) Actividades clave → (8) Socios clave → (9) Estructura de costos. Empezar por el segmento es regla de oro: si el segmento es difuso, todo el Canvas será difuso. Los novatos suelen empezar por la propuesta de valor (\emph{``mi idea es…''}) y por eso construyen para un cliente que no existe.}
\milcpaso{3}{milcTurquesa}{Lo esencial cabe en una pregunta: \emph{¿hay coherencia entre los 9 bloques?}. Coherencia significa que la propuesta de valor encaja con el dolor del segmento, los canales son donde el segmento está, los recursos clave alcanzan para sostener las actividades clave, y los ingresos cubren los costos. Un Canvas coherente se lee como una historia: \emph{para [segmento], que sufre [dolor], ofrecemos [propuesta] que llega por [canales], cobramos [ingresos] y nos cuesta [costos]}.}
\milcpaso{4}{milcTurquesa}{Algoritmo: (1) dibuja el Canvas en una hoja apaisada con 9 cuadros; (2) llena Segmento con 3-4 rasgos; (3) escribe la Propuesta de valor en 1 frase \emph{``ayudamos a X a Y para que Z''}; (4) anota 2-3 Canales reales (no \emph{``redes sociales''}: TikTok, WhatsApp Business, voz a voz en la plaza); (5) define 1 tipo de Relación (transaccional, suscripción, atención personal); (6) define Fuente de ingreso con número aproximado; (7) lista Recursos, Actividades y Socios; (8) cierra con Costos; (9) léelo completo y verifica coherencia.}

\begin{drawbox}{milcTurquesa}{Anatomía del Canvas profesional}
\textbf{Segmento:} ¿quién es el cliente exacto? \\[1mm] \textbf{Propuesta de valor:} ¿qué dolor le quitas? \\[1mm] \textbf{Canales:} ¿cómo lo alcanzas? \\[1mm] \textbf{Relación:} ¿cómo interactúan? \\[1mm] \textbf{Ingresos:} ¿cómo y cuánto te paga? \\[1mm] \textbf{Recursos clave:} ¿qué necesitas tener? \\[1mm] \textbf{Actividades clave:} ¿qué haces tú? \\[1mm] \textbf{Socios clave:} ¿quién hace lo que tú no? \\[1mm] \textbf{Costos:} ¿qué te cuesta entregar?
\end{drawbox}

\begin{softbox}{milcMostaza}{milcGris}{Errores comunes y su corrección}
(1) Empezar por la propuesta de valor antes que por el segmento: produce Canvas para clientes ficticios. (2) Segmento genérico (\emph{``los jóvenes''}, \emph{``las empresas''}): inservible para tomar decisiones. (3) Propuesta de valor sin verbo de cambio (\emph{``una app de…''} no es propuesta de valor; \emph{``ahorra 4 horas semanales a los tenderos''} sí lo es). (4) Canales inventados sin evidencia (poner \emph{``redes sociales''} cuando ningún entrevistado mencionó usar Instagram). (5) Ingresos sin número (decir \emph{``cobraremos''} sin saber cuánto). (6) Costos olvidados o subestimados (no incluir el tiempo del propio emprendedor). (7) Canvas hecho de un solo intento, sin revisarlo en voz alta para detectar incoherencias.
\end{softbox}

\begin{infoband}{milcTurquesa}


\textbf{Lo que va al cuaderno (Actividad 2 · EXPLICA).} Título: «Actividad 2 --- Anatomía del Canvas». \textbf{Formato:} (a) lista de los 9 bloques con 1 pregunta clave por bloque; (b) orden profesional de llenado numerado; (c) los 7 errores comunes con marca del que más temes cometer. \textbf{Sabes que terminaste cuando} podrías explicarle el Canvas a un familiar sin abrir esta guía.
\end{infoband}

\puente{Tienes el método. Toca dibujar el Canvas en una hoja y llenar los 9 bloques con frases concretas, no genéricas. Termina identificando los 3 supuestos más críticos: los bloques cuya falsedad rompería todo el modelo y que la sesión 4 (prototipo) tendrá que verificar primero.}

\milcestacion{milcMagenta}{Estación 3 de 3 · Hacer}

\begin{softbox}{milcMagenta}{milcGris}{Cómo vas a construir tu producto}
\textbf{Actividad 3 · CREA + EVALÚA --- Mi Canvas + 3 supuestos críticos} (35 min · individual). Hora de construir. Dibujas tu Canvas en una hoja apaisada, llenas los 9 bloques en orden profesional, lo lees completo en voz alta para verificar coherencia y cierras identificando los 3 supuestos más críticos: aquellos bloques cuya falsedad haría caer todo el modelo.
\end{softbox}

\milcpaso{1}{milcMagenta}{Tu Canvas se construye en 5 pasos: (1) dibujar en hoja apaisada los 9 cuadros con sus nombres; (2) llenar en el orden profesional (Segmento → ... → Costos); (3) usar frases cortas y concretas (3-7 palabras por idea, con post-its mentales); (4) anclar cada bloque a evidencia de las entrevistas cuando sea posible; (5) leer el Canvas completo como historia para verificar coherencia.}
\milcpaso{2}{milcMagenta}{Sigue el patrón \textbf{``frase corta sobre párrafo largo''}: los Canvas profesionales se leen en 2 minutos. Si un bloque tiene un párrafo, está mal hecho. Cada bloque debería tener 3-5 ideas en frases cortas, separadas. Esa disciplina visual obliga a la precisión y revela cuando un bloque es relleno.}
\milcpaso{3}{milcMagenta}{Los 3 supuestos críticos se identifican preguntando: \emph{``si este bloque resulta falso, ¿cuánto se cae el modelo?''}. Los bloques que más suelen contener supuestos críticos son: (a) \textbf{Propuesta de valor}: supuesto de que el cliente quiere lo que tú ofreces; (b) \textbf{Canales}: supuesto de que sabes dónde encontrarlo; (c) \textbf{Fuentes de ingreso}: supuesto de que pagará lo que pides. Estos 3 son los blancos preferidos del prototipo de la sesión 4.}
\milcpaso{4}{milcMagenta}{Algoritmo: (1) terminado el Canvas, marca con asterisco los 3 bloques más débiles (los que llenaste con menos evidencia); (2) redacta cada supuesto crítico como afirmación falsable: \emph{``los tenderos pagarán 20 mil pesos al mes por este servicio''}; (3) anota cómo verificarías cada supuesto en 1 semana (oferta directa, demo pagada, intención escrita); (4) cierra con la frase: \emph{``si los 3 supuestos críticos sobreviven, mi modelo es sólido''}.}

\begin{drawbox}{milcMagenta}{Checklist antes de entregar el Canvas}
\checkbox Canvas dibujado en hoja apaisada con los 9 bloques visibles. \\[1mm] \checkbox Bloques llenos en el orden profesional (Segmento primero). \\[1mm] \checkbox Frases cortas (3-7 palabras), no párrafos. \\[1mm] \checkbox Cada bloque anclado a evidencia de las entrevistas cuando aplica. \\[1mm] \checkbox Coherencia verificada leyendo el Canvas en voz alta. \\[1mm] \checkbox 3 supuestos críticos identificados con asterisco. \\[1mm] \checkbox Tabla con los 3 supuestos formulados como afirmaciones falsables + cómo verificarlos.
\end{drawbox}

\begin{softbox}{milcMagenta}{milcGris}{Plantilla de trabajo}
\textbf{Segmento.} \textit{3-4 rasgos concretos.} \\[1mm] \textbf{Propuesta de valor.} \textit{Ayudamos a X a Y para que Z.} \\[1mm] \textbf{Canales.} \textit{2-3 canales reales.} \\[1mm] \textbf{Relación.} \textit{Transaccional / suscripción / atención personal.} \\[1mm] \textbf{Ingresos.} \textit{Cómo cobras + monto aproximado.} \\[1mm] \textbf{Recursos clave.} \textit{Lista corta.} \\[1mm] \textbf{Actividades clave.} \textit{Lista corta.} \\[1mm] \textbf{Socios clave.} \textit{Lista corta.} \\[1mm] \textbf{Costos.} \textit{Lista corta con números.}
\end{softbox}

\begin{infoband}{milcMagenta}


\textbf{Lo que va al cuaderno (Actividad 3 · CREA + EVALÚA).} Título: «Actividad 3 --- Mi Canvas + supuestos críticos». \textbf{Formato:} (a) Canvas dibujado en hoja apaisada (pegada o adjunta al cuaderno); (b) tabla 3×3 con los 3 supuestos críticos (Supuesto~|~Bloque del Canvas~|~Cómo verificarlo en 1 semana). \textbf{Sabes que terminaste cuando} podrías presentarle el Canvas a tu familia en 2 minutos y entender qué vas a probar primero.
\end{infoband}

\puente{Lo que sigue resume \emph{qué} entregas hoy: Canvas en una hoja con los 9 bloques llenos + tabla de 3 supuestos críticos. El criterio principal: que un asesor financiero leyendo tu Canvas pueda decir en 3 minutos \emph{``entendí qué vendes, a quién, cómo se entera, cómo te paga y cuánto te cuesta''}.}

\milcestacion{milcMostaza}{Producto de la sesión}
\textbf{Canvas con 9 bloques llenos + 3 supuestos críticos formulados como afirmaciones falsables}.
\begin{softbox}{milcMostaza}{milcGris}{Criterios de calidad}
(1) 9 bloques llenos con frases cortas y concretas. (2) Orden profesional de llenado seguido (Segmento primero). (3) Bloques anclados a evidencia de las entrevistas. (4) Coherencia verificada en voz alta. (5) 3 supuestos críticos identificados. (6) Plan de verificación de 1 semana para cada supuesto.
\end{softbox}

\puente{Antes de cerrar, mira el modelo desde las cinco dimensiones humanas. Diseñar el modelo con respeto al cliente es ética del oficio; diseñarlo extrayendo valor sin devolverlo es la peor versión del emprendimiento contemporáneo.}

\milcestacion{milcVino}{Evaluación liberadora · cinco dimensiones}

\begin{softbox}{milcVino}{milcGris}{Sentido de la evaluación}
Evaluar no es solo revisar una nota. Es reconocer qué aprendí, qué cambió en mí, qué emociones manejé, cómo me relaciono con la ciudad y la comunidad, y qué le dejaré a quien viene detrás.
\end{softbox}

\textbf{Mi reflexión liberadora en cinco dimensiones.} Completa cada frase iniciada:

\small
\begin{tabularx}{\textwidth}{>{\bfseries\RaggedRight\arraybackslash}p{3.4cm}Y>{\RaggedRight\arraybackslash}p{4.6cm}}
\rowcolor{milcVino}\color{white}Dimensión & \color{white}\bfseries Mi respuesta (frame starter) & \color{white}\bfseries Algo que descubrí\\
1. Personal & Lo que más cambió en mí fue \linead{6.5cm} & \linead{4.2cm}\\
2. Emocional & La emoción más fuerte fue \linead{2.5cm} y la regulé \linead{3cm} & \linead{4.2cm}\\
3. Ciudadana & En democracia esto importa porque \linead{6cm} & \linead{4.2cm}\\
4. Local (Cartago) & En mi barrio sería útil para \linead{6.5cm} & \linead{4.2cm}\\
5. Intergeneracional & A mi abuela/abuelo le diría \linead{6.5cm} & \linead{4.2cm}\\
\end{tabularx}
\normalsize

\begin{infoband}{milcVino}
\textbf{Para la próxima sustentación.} Vuelve a esta tabla en seis meses. Verás que las respuestas cambian — no porque lo aprendido cambie, sino porque tú cambias. La evaluación liberadora es una conversación contigo a través del tiempo.
\end{infoband}


\puente{Y para terminar, tres voces sobre el valor y la economía. Dussel sobre los modelos que liberan o extraen, Aristóteles sobre la phronesis económica (el justo medio entre lucro y servicio), Floridi sobre la transparencia del modelo como ética profesional.}

\milcestacion{milcGradeDark}{Triángulo de pensamiento · cierre filosófico}

\begin{softbox}{milcVino}{milcGris}{Dussel · lente del nosotros}
\textit{Todo modelo de negocio decide si libera al cliente o lo extrae; no hay diseño económico neutro.}

{\footnotesize\itshape\color{milcNegro!70}--- formulación de esta guía, al modo de Enrique Dussel. No es una cita textual.}

Tu Canvas es una toma de posición ética antes que comercial. Hay modelos que liberan (te dan valor real por un pago justo) y modelos que extraen (capturan tu atención, tu tiempo o tu dinero sin devolver proporción). La filosofía de la liberación pide que el Canvas declare con honestidad de qué lado estás: ¿la propuesta de valor mejora la vida del segmento o explota su vulnerabilidad? La phronesis del emprendedor consiste en construir modelos que devuelven valor proporcional a lo que extraen.

\textbf{Mi Canvas, ¿libera o extrae? ¿La propuesta de valor mejora la vida del segmento o aprovecha su falta de opciones?}
\end{softbox}
\linead{\linewidth}\\[1mm]
\linead{\linewidth}

\begin{softbox}{milcMostaza}{milcGris}{Estoicismo · Séneca · cuidado interior}
\textit{No hay viento favorable para quien no sabe a qué puerto va; pero tampoco hay sostenibilidad para quien navega ignorando los costos.}

{\footnotesize\itshape\color{milcNegro!70}--- formulación de esta guía, al modo de Séneca. No es una cita textual.}

Es tentador inflar la fuente de ingresos para que el Canvas \emph{cierre} con números bonitos, o ignorar los costos para que el modelo parezca rentable. La sabiduría estoica recomienda el justo medio: ni un modelo que solo busca lucro (extracción sin servicio), ni un modelo que solo sirve sin sostenerse (servicio sin futuro). Tu Canvas debería poder sostenerse 3 años sin sacrificar al cliente ni al emprendedor.

\textbf{¿Mi Canvas se sostiene económicamente sin volverse extractivo? ¿Estoy cobrando justo o estoy regalando trabajo?}
\end{softbox}
\linead{\linewidth}\\[1mm]
\linead{\linewidth}

\begin{softbox}{milcTurquesa}{milcGris}{Floridi · lente de la infoesfera}
\textit{La transparencia del modelo es la nueva forma de la ética profesional contemporánea.}

{\footnotesize\itshape\color{milcNegro!70}--- formulación de esta guía, al modo de Luciano Floridi. No es una cita textual.}

Un Canvas que el cliente nunca podría ver sin sentirse engañado no es un modelo sostenible: es un modelo opaco. La ética digital contemporánea exige que el emprendedor pueda mostrar su Canvas al cliente y este lo apruebe sin escándalo. Esa transparencia es virtud técnica y comercial. Los modelos opacos sobreviven a corto plazo y se derrumban cuando la información se vuelve pública.

\textbf{Si mi cliente viera mi Canvas completo (con costos, márgenes y socios), ¿se sentiría respetado o engañado?}
\end{softbox}
\linead{\linewidth}\\[1mm]
\linead{\linewidth}

\begin{infoband}{milcNegro}
\textbf{Nota para el docente.} Las tres formulaciones son del autor de la guía, no citas textuales de los pensadores: recogen su lente para pensar el tema, no sus palabras. Si vas a citarlos en clase, remite a la obra original.
\end{infoband}

\milcestacion{milcVino}{Mi autoevaluación · marca una casilla por fila}

{\small
\setlength{\extrarowheight}{2.2pt}
\begin{tabularx}{\textwidth}{>{\RaggedRight\arraybackslash\bfseries}p{3.0cm}YYY}
\rowcolor{milcVino}\color{white}\bfseries Criterio & \color{white}\bfseries Lo logré & \color{white}\bfseries En proceso & \color{white}\bfseries Necesito apoyo\\[.6mm]
Comprensión del tema & \milcopcion{Explico las ideas con mis palabras.} & \milcopcion{Reconozco las ideas, falta precisión.} & \milcopcion{Necesito repasar lo básico.}\\[.8mm]
\rowcolor{milcGris}Pensamiento computacional & \milcopcion{Usé los pilares al armar mi producto.} & \milcopcion{Usé algunos pilares.} & \milcopcion{Necesito releer la estación 2.}\\[.8mm]
Aplicación y producto & \milcopcion{Mi producto cumple los criterios.} & \milcopcion{Está armado, con ajustes pendientes.} & \milcopcion{Aún está en construcción.}\\[.8mm]
\rowcolor{milcGris}Reflexión 5 dimensiones & \milcopcion{Respondí cada una con honestidad.} & \milcopcion{Respondí algunas con detalle.} & \milcopcion{Mis respuestas son muy generales.}\\[.8mm]
Triángulo de pensamiento & \milcopcion{Conecté las 3 voces con mi trabajo.} & \milcopcion{Conecté al menos una.} & \milcopcion{No las relaciono con mi proceso.}\\[.8mm]
\end{tabularx}}

\vspace{3mm}
\textbf{Hoy entendí que:}\\[1mm]
\linead{\linewidth}\\[1mm]
\linead{\linewidth}

\textbf{Mi compromiso para la próxima sesión es:}\\[1mm]
\linead{\linewidth}\\[1mm]
\linead{\linewidth}

\begin{softbox}{milcMostaza}{milcGris}{Cita de despedida · Marco Aurelio}
\textit{``El obstáculo en el camino se vuelve el camino. Lo que estorba al camino se vuelve camino.''}\\[1mm]
Cada error de hoy, bien observado, se convierte en pieza del camino que pisarás mañana.
\end{softbox}

\small
\begin{tabularx}{\textwidth}{>{\bfseries}p{3.6cm}Y}
\rowcolor{milcGradeDark!12}Nombre completo: & \linead{10cm}\\
Fecha: & \linead{4cm} \quad Curso/grupo: \linead{4cm}\\
Mi firma: & \linead{9cm}\\
\end{tabularx}
\normalsize

\begin{infoband}{milcGradeDark}
\textbf{Para la próxima sesión.} Llega con el Canvas terminado y los 3 supuestos críticos formulados. En la sesión 4 vas a construir el \textbf{prototipo mínimo viable} que verifique uno o varios de esos supuestos. El Canvas se honra construyendo lo que prueba sus afirmaciones, no archivándolo en una carpeta.
\end{infoband}

\end{document}
